Consider now a jump-diffusion process
\begin{equation} \label{JumpDiffusionProcess}
	X_t = \sigma W_t + \mu t + J_t = X^c(t) + J_t,
\end{equation}
where $J$ is a compound Poisson process and $X^c$ is the continuous part of $X$:
\begin{equation} \label{ComponentDefinitions}
	J_t = \sum_{i=1}^{N_t} \Delta X_i, \qquad\qquad X^c_t = \mu t + \sigma W_t.
\end{equation}
Define $Y_t = f(X_t)$ where $f \in C^2(\mathbb{R})$ and denote by $T_i, i = 1 \dots N_T$ the jump times of $X$. On $]T_i, T_{i+1}[$, $X$ evolves according to
\begin{equation} \label{EvolutionBetweenJumps}
	dX_t = dX^c_t = \sigma dW_t + \mu dt,
\end{equation}
so by applying the \Ito{} formula in the Brownian case we obtain
\begin{align} \label{ItoIntegration}
    Y_{T_{i+1}-} - Y_{T_i} &= \int_{T_i}^{T_{i+1}-} \frac{\sigma^2}{2} f''(X_t) \di t + \int_{T_i}^{T_{i+1}-} f'(X_t) \di X_t \nonumber \\
    &= \int_{T_i}^{T_{i+1}-} \left\{ \frac{\sigma^2}{2} f''(X_t) \di t + f'(X_t) \di X^c_t \right\}
\end{align}
since $dX_t = dX^c(t)$ on this interval. If a jump of size $\Delta X_t$ occurs then the resulting change in $Y_t$ is given by $f(X_{t-} + \Delta X_t) - f(X_{t-})$. The total change in $Y_t$ can therefore be written as the sum of these two contributions:
\begin{equation} \label{FinalItoFormula}
	\begin{split}
		f(X_t) - f(X_0) &= \int_0^t f'(X_s) dX^c_s + \int_0^t \frac{\sigma^2}{2} f''(X_s) ds \\
		&\quad + \sum_{0 \le s \le t, \Delta X_s \neq 0} [f(X_{s-} + \Delta X_s) - f(X_{s-})].
	\end{split}
\end{equation}

\begin{xattention}{}
Replacing $\di X_s^c$ by $\di X_s - \Delta X_s$ one obtains an equivalent expression:
\begin{equation} \label{EquivalentItoExpression}
	\begin{split}
		f(X_t) - f(X_0) &= \int_0^t f'(X_{s-}) \di X_s + \int_0^t \frac{\sigma^2}{2} f''(X_s) \di s \\
		&\quad + \sum_{0 \le s \le t, \Delta X_s \neq 0} [f(X_{s-} + \Delta X_s) - f(X_{s-}) - \Delta X_s f'(X_{s-})].
	\end{split}
\end{equation}
When the number of jumps is finite, this form is equivalent to \cref{FinalItoFormula}. However, as we will see below, the form \cref{EquivalentItoExpression} is more general, wherein both the stochastic integral and the sum over the jumps in \cref{EquivalentItoExpression} are well-defined for any semi-martingale, even if the number of jumps is infinite, while the sum in \cref{FinalItoFormula} may not converge if jumps have infinite variation.

Here we have only used the \Ito{} formula for diffusions, which is of course still valid if $\sigma$ is replaced by a process $(\sigma_t)_{t \in [0,T]}$ belonging to the class $\mathcal{P}_T$.
\end{xattention}

\begin{prop}{\Ito{} Formula for Jump-Diffusion Processes}{ItoJumpDiffusion}
Let $X$ be a diffusion process with jumps, defined as the sum of a drift term, a Brownian stochastic integral and a compound Poisson process:
\begin{equation*}
	X_t = X_0 + \int_0^t b_s \di s + \int_0^t \sigma_s \di W_s + \sum_{i=1}^{N_t} \Delta X_i,
\end{equation*}
where $b_t$ and $\sigma_t$ are continuous adapted processes with
\begin{equation*}
	\E\left[ \int_0^T \sigma_t^2 \di t \right] < \infty.
\end{equation*}
Then, for any $C^{1,2}$ function $f : [0,T] \times \R \to \R$, the process $Y_t = f(t, X_t)$ can be represented as
\begin{align*}
	f(t, X_t) - f(0, X_0) &= \int_0^t \left[ \frac{\partial f}{\partial s}(s, X_s) + \frac{\partial f}{\partial x}(s, X_s)b_s \right] \di s \\
	&\quad + \frac{1}{2} \int_0^t \sigma_s^2 \frac{\partial^2 f}{\partial x^2}(s, X_s)\di s + \int_0^t \frac{\partial f}{\partial x}(s, X_s)\sigma_s \di W_s \\
	&\quad + \sum_{\{i \ge 1, T_i \le t\}} [f(X_{T_i-} + \Delta X_i) - f(X_{T_i-})].
\end{align*}
Expressed in differential notation
\begin{align*}
	\di Y_t &= \frac{\partial f}{\partial t}(t, X_t) \di t + b_t \frac{\partial f}{\partial x}(t, X_t) \di t + \frac{\sigma_t^2}{2} \frac{\partial^2 f}{\partial x^2}(t, X_t) \di t \\
	&\quad + \frac{\partial f}{\partial x}(t, X_t) \sigma_t \di W_t + [f(X_{t-} + \Delta X_t) - f(X_{t-})].
\end{align*}
\end{prop}

\begin{defn}{Infinite Divisibility}{InfiniteDivisibility}
    A probability distribution $F$ on $\R^d$ is said to be \textit{infinitely divisible} if for any integer $n \ge 2$, there exist $n$ $\iid$ random variables $Y_1, \dots, Y_n$ such that $Y_1 + \dots + Y_n$ has distribution $F$.
\end{defn}

\begin{prop}{Infinite Divisibility and \Levy{} Processes}{InfiniteDivisibilityAndLevyProcesses}
    Let $X \coloneqq (X_t)_{t \ge 0}$ be a  \Levy{} process. Then for every $t$, $X_t$ has an infinitely divisible distribution. Conversely, if $F$ is an infinitely divisible distribution, there exists a $\Levy$ process $X$ such that the distribution of $X_1$ is given by $F$.
\end{prop}

\begin{prop}{Characteristic function of a \Levy{} process}{CharacteristicFunctionLevy}
Let $(X_t)_{t\geq 0}$ be a \Levy{} process on $\R^d$. There exists a continuous function $\psi : \R^d \mapsto \R$ called the characteristic exponent of $X$, such that:
\begin{equation}\label{eq:CharacteristicFunction}
    E\left[e^{iz \cdot X_t}\right] = e^{t\psi(z)}, \quad z \in \R^d.
\end{equation}
\end{prop}

\section{Integro-differential equations and numerical methods}
\epigraph{An approximate answer to the right problem is worth a good deal more than an exact answer to an approximate problem.}{John Tukey}
